% Source: Wald GR, physical PDF pp. 104--108 (printed pp. 91--95).
% Coverage: Section 5.1, Homogeneity and Isotropy; equations (5.1.1)--(5.1.11).
% Figures/tables/footnotes: Figures 5.1--5.2.
% Status: source-reviewed and compile-clean.
% Uncertainties: none.

\subsection{Homogeneity and Isotropy}\label{sec:5.1}

In the subject of cosmology, it is very difficult to prove theories by appealing only
to observational data. We have direct contact in our lifetime or even in the lifetime
of human civilization with only a negligibly small spacetime region of our universe.
While our telescopes can observe objects remarkably far away by ordinary human
scales, it should be recognized that in cosmic terms they report information about
only a portion of our past light cone. Thus, a good deal of our input in the subject
of cosmology arises from our philosophical prejudices. Observational data may
confirm these prejudices, but in general they cannot be expected to definitively prove
that they are correct. Nevertheless, the cosmological models considered in this
chapter have provided a remarkably successful account of the nature of the universe.

Since the time of Copernicus, it has generally been assumed that we do not occupy
a privileged position in our universe; that if we were located in a different region of
our universe, the basic characteristics of our surroundings would appear the same.
Similarly, it is natural to assume that the universe is isotropic, that is, that there are
no preferred directions in space; that observations on sufficiently large scales should
yield results which do not depend on which direction we look. These philosophical
prejudices of homogeneity and isotropy have received strong confirmation from
modern observations. While observations of the distribution of galaxies in our
universe show clustering of galaxies on a wide range of distance scales and recent
observations have shown large regions devoid of galaxies (Kirshner et al. 1981), on
the largest scales the galaxy distribution appears to be homogeneous and isotropic.
Counts of radio sources and the isotropy of the X-ray and \(\gamma\)-ray background
radiation also support the hypothesis of homogeneity and isotropy of the universe on
large scales. Even stronger observational evidence for the homogeneity and isotropy
of our universe comes from the discovery of thermal radiation at about \(3\,\mathrm{K}\)
filling our universe, which has been measured to be isotropic to a very high
precision. As discussed in Section~5.4, this radiation is believed to have a
cosmological origin, and it would be very difficult to explain its existence and its
isotropy if the hypothesis of the homogeneity and isotropy of the universe were not
valid to a very good approximation on large distance scales.

Thus, for the remainder of this chapter, we shall proceed under the assumption that
the universe is homogeneous and isotropic. Our first task is to formulate precisely
the mathematical meaning of this assumption. Loosely speaking, homogeneity means
that at any given ``instant of time'' each point of ``space'' should ``look like'' any
other point. A precise formulation can be given as follows: A spacetime is said to be
(spatially) homogeneous if there exists a one-parameter family of spacelike
hypersurfaces \(\Sigma_t\) foliating the spacetime (see Figure~\ref{fig:5.1}) such that
for each \(t\) and for any points \(p,q\in\Sigma_t\) there exists an isometry of the
spacetime metric, \(g_{ab}\), which takes \(p\) into \(q\). (See Appendix~C for the
definition of an isometry.)

% Source: Wald GR, physical PDF p. 105
%         (printed p. 92).
% Coverage: Figure 5.1.
% Figures/tables/footnotes: Figure 5.1; no tables or footnotes.
% Status: source-reviewed and compile-clean.
% Uncertainties: none.
\begingroup
\renewcommand{\thefigure}{5.1}
\begin{figure}[htbp]
  \centering
  \begin{tikzpicture}[>=Stealth, line cap=round, line join=round, scale=.92]
    \draw[thick] (-3.0,0.15) .. controls (-1.7,-.55) and (1.0,-.55) .. (3.0,0.15)
      .. controls (1.4,.75) and (-1.6,.75) .. cycle;
    \draw[thick] (-2.55,1.35) .. controls (-1.25,.73) and (1.25,.73) .. (2.55,1.35)
      .. controls (1.18,1.96) and (-1.22,1.96) .. cycle;
    \draw[thick] (-2.05,2.55) .. controls (-.95,1.98) and (.95,1.98) .. (2.05,2.55)
      .. controls (.98,3.08) and (-1.00,3.08) .. cycle;
    % Representative isometries act within each homogeneous slice.  Use
    % several short copies, as in the source, rather than a single trajectory
    % that could be mistaken for a preferred flow.
    \foreach \xa/\xb in {-2.15/-.95,-.65/.55,.85/2.05} {
      \draw[->, thick] (\xa,.12) .. controls ({(\xa+\xb)/2},.39) .. (\xb,.15);
    }
    \foreach \xa/\xb in {-1.48/-.58,-.42/.48,.64/1.54} {
      \draw[->, thick] (\xa,2.52) .. controls ({(\xa+\xb)/2},2.75) .. (\xb,2.54);
    }
    \fill (-.75,1.34) circle (1.4pt) node[below left=2pt] {\(p\)};
    \fill (1.10,1.34) circle (1.4pt) node[below right=2pt] {\(q\)};
    \draw[->, thick] (-.65,1.37) .. controls (0,1.62) and (.52,1.60) .. (1.00,1.37);
    \draw[->, thick] (-2.05,1.35) .. controls (-1.72,1.56) and (-1.38,1.56) .. (-1.08,1.37);
    \draw[->, thick] (1.36,1.37) .. controls (1.68,1.57) and (1.98,1.55) .. (2.27,1.37);
    \node[right] at (2.95,1.60) {\(\Sigma_t\)};
  \end{tikzpicture}
  \caption{The hypersurfaces of spatial homogeneity in spacetime. By definition of
  homogeneity, for each \(t\) and each \(p,q\in\Sigma_t\), there exists an isometry of
  the spacetime which takes \(p\) into \(q\).}
  \label{fig:5.1}
\end{figure}
\endgroup


With regard to isotropy, it first should be pointed out that, in general, at each
point, at most one observer can see the universe as isotropic. For example, if
ordinary matter fills the universe, any observer in motion relative to the matter must
see an anisotropic velocity distribution of the matter. With this fact in mind, a precise
formulation of the notion of isotropy can be given as follows: A spacetime is said to
be (spatially) isotropic at each point if there exists a congruence of timelike curves
(i.e., observers), with tangents denoted \(u^a\), filling the spacetime (see
Figure~\ref{fig:5.2}) and satisfying the following property. Given any point \(p\) and
any two unit ``spatial'' tangent vectors \(s_1^a,s_2^a\in T_pM\) (i.e., vectors at
\(p\) orthogonal to \(u^a\)), there exists an isometry of \(g_{ab}\) which leaves \(p\)
and \(u^a\) at \(p\) fixed but rotates \(s_1^a\) into \(s_2^a\). Thus, in an isotropic
universe it is impossible to construct a geometrically preferred tangent vector
orthogonal to \(u^a\).

% Source: Wald GR, physical PDF p. 106
%         (printed p. 93).
% Coverage: Figure 5.2.
% Figures/tables/footnotes: Figure 5.2; no tables or footnotes.
% Status: source-reviewed and compile-clean.
% Uncertainties: none.
\begingroup
\renewcommand{\thefigure}{5.2}
\begin{figure}[htbp]
  \centering
  \begin{tikzpicture}[>=Stealth, line cap=round, line join=round, scale=.90]
    \foreach \x/\bend in {-2.3/-16,-1.55/-10,-.78/-5,0/0,.78/5,1.55/10,2.3/16} {
      \draw[thick] (\x,-1.30) .. controls (\x+\bend/35,.15) and
        (\x-\bend/35,1.20) .. (\x,2.85);
    }
    \fill (0,.62) circle (1.5pt) node[below left=2pt] {\(p\)};
    \draw[->, thick] (0,.62) -- (.05,1.52) node[left] {\(u^a\)};
    \draw[->, thick] (0,.62) -- (.62,1.08) node[above right] {\(s_1^a\)};
    \draw[->, thick] (0,.62) -- (1.02,.62) node[below right] {\(s_2^a\)};
    \draw[->, dashed] (.55,1.01) to[bend left=25] (.88,.67);
  \end{tikzpicture}
  \caption{The world lines of isotropic observers in spacetime. By definition of
  isotropy, for any two vectors \(s_1^a,s_2^a\) at \(p\) which are orthogonal to
  \(u^a\), there exists an isometry of the spacetime which leaves \(p\) fixed and
  rotates \(s_1^a\) into \(s_2^a\).}
  \label{fig:5.2}
\end{figure}
\endgroup


It is not difficult to see that in the case of a homogeneous and isotropic spacetime,
the surfaces \(\Sigma_t\) of homogeneity must be orthogonal to the tangents,
\(u^a\), to the world lines of the isotropic observers. If not, then assuming that the
isotropic observers and the family of homogeneous surfaces \(\Sigma_t\) are unique,
the failure of the tangent subspace orthogonal to \(u^a\) to coincide with the tangent
space of \(\Sigma_t\) would enable us to construct a geometrically preferred spatial
vector, in violation of isotropy. (If the isotropic observers or family of homogeneous
surfaces are not unique, as in special cases such as flat spacetime, it is still possible
to show the existence of a family of isotropic observers orthogonal to a family of
homogeneous surfaces.) Now, the spacetime metric, \(g_{ab}\), induces a Riemannian
metric, \(h_{ab}(t)\), on each \(\Sigma_t\) by restricting the action of \(g_{ab}\) at each
\(p\in\Sigma_t\) to vectors tangent to \(\Sigma_t\). The induced spatial geometry of
the surfaces \(\Sigma_t\) is greatly restricted by the following requirements: (i)
Because of homogeneity, there must be isometries of \(h_{ab}\) which carry any
\(p\in\Sigma_t\) into any \(q\in\Sigma_t\). (ii) Because of isotropy, it must be
impossible to construct any geometrically preferred vectors on \(\Sigma_t\).

We shall now show that the second requirement implied by isotropy is particularly
restrictive. Consider the Riemann tensor \({}^{(3)}R_{abcd}\) constructed from
\(h_{ab}\) on \(\Sigma_t\). If we raise the third index with \(h_{ab}\), we may view
\({}^{(3)}R_{ab}{}^{cd}\) at point \(p\) as a linear map, \(L\), of the vector space
\(W\) of two-forms [i.e., antisymmetric tensors of rank \((0,2)\)] at \(p\) into
itself \(L:W\to W\). By Eq.~\eqref{eq:3.2.20}, \(L\) is symmetric, i.e., it is a
self-adjoint map (with the natural, positive definite inner product on \(W\)
determined by \(h_{ab}\)). Therefore, \(W\) has an orthonormal basis of eigenvectors
of \(L\). If the eigenvalues of these eigenvectors were distinct, then we would be
able to give a geometrical prescription for picking out a preferred two-form at \(p\)
and, consequently, a preferred vector at \(p\). Hence, in order not to violate
isotropy, all the eigenvalues of \(L\) must be equal. This means that \(L\) is a
multiple of the identity operator,
\begin{equation}
  L=KI.
  \tag{5.1.1}\label{eq:5.1.1}
\end{equation}
That is,
\begin{equation}
  {}^{(3)}R_{ab}{}^{cd}=K\delta^c_{[a}\delta^d_{b]}.
  \tag{5.1.2}\label{eq:5.1.2}
\end{equation}
Thus, lowering the indices, we have
\begin{equation}
  {}^{(3)}R_{abcd}=K h_{c[a}h_{b]d}.
  \tag{5.1.3}\label{eq:5.1.3}
\end{equation}
The requirement (i) of homogeneity implies that \(K\) must be a constant, i.e., it
cannot vary from point to point of \(\Sigma_t\). Actually, it is an interesting fact
that the requirement (ii) of isotropy at each point also implies the constancy of
\(K\). To prove this, we substitute Eq.~\eqref{eq:5.1.3} into the Bianchi identity,
Eq.~\eqref{eq:3.2.16}, to obtain,
\begin{equation}
  0=D_{[e}{}^{(3)}R_{ab]cd}=(D_{[e}K)h_{|c|a}h_{b]d},
  \tag{5.1.4}\label{eq:5.1.4}
\end{equation}
where \(D_a\) denotes the derivative operator on \(\Sigma_t\) associated with
\(h_{ab}\). (We use the notation \(D_a\) rather than \(\nabla_a\) in order to avoid
confusion with the derivative operator on the four-dimensional spacetime associated
with \(g_{ab}\).) On a manifold of dimension three or greater, the right side of
Eq.~\eqref{eq:5.1.4} will vanish if and only if \(D_aK=0\), i.e., \(K\) is constant.
Thus, we can actually dispense with the assumption of homogeneity in discussing the
geometry of \(\Sigma_t\).

A space where Eq.~\eqref{eq:5.1.3} is satisfied (with \(K=\text{constant}\)) is
called a space of constant curvature. It can be shown (Eisenhart 1949) that any two
spaces of constant curvature of the same dimension and metric signature which have
equal values of \(K\) must be (locally) isometric. Thus, our task of determining the
possible spatial geometries of \(\Sigma_t\) will be completed if we enumerate spaces
of constant curvature encompassing all values of \(K\). This is easily done. All
positive values of \(K\) are attained by the 3-spheres, defined as the surfaces in
four-dimensional flat Euclidean space \(\mathbb{R}^4\) whose Cartesian coordinates
satisfy
\begin{equation}
  x^2+y^2+z^2+w^2=R^2.
  \tag{5.1.5}\label{eq:5.1.5}
\end{equation}
In spherical coordinates, the metric of the unit 3-sphere is
\begin{equation}
  ds^2=d\psi^2+\sin^2\psi\bigl(d\theta^2+\sin^2\theta\,d\phi^2\bigr).
  \tag{5.1.6}\label{eq:5.1.6}
\end{equation}
The value \(K=0\) is attained by ordinary three-dimensional flat space. In Cartesian
coordinates, this metric is
\begin{equation}
  ds^2=dx^2+dy^2+dz^2.
  \tag{5.1.7}\label{eq:5.1.7}
\end{equation}

Finally, all negative values of \(K\) are attained by the three-dimensional
hyperboloids, defined as the surfaces in a four-dimensional flat Lorentz signature
space (i.e., Minkowski spacetime) whose global inertial coordinates satisfy
\begin{equation}
  t^2-x^2-y^2-z^2=R^2.
  \tag{5.1.8}\label{eq:5.1.8}
\end{equation}
In hyperbolic coordinates, the metric of the unit hyperboloid is
\begin{equation}
  ds^2=d\psi^2+\sinh^2\psi\bigl(d\theta^2+\sin^2\theta\,d\phi^2\bigr).
  \tag{5.1.9}\label{eq:5.1.9}
\end{equation}

The new possibilities for the global spatial structure of our universe should be
stressed. In prerelativity physics, as well as in special relativity, it was assumed
that space had the flat structure given by the possibility \(K=0\) above. But even
under the very restrictive assumptions of homogeneity and isotropy, the framework
of general relativity admits two other distinct possibilities. The possibility of a
3-sphere spatial geometry is particularly interesting, as it is a compact manifold
(see Appendix~A) and thus describes a universe which is finite but has no boundary.
Such a universe is called ``closed,'' while the universes with noncompact spatial
sections such as those given by flat and hyperboloid geometries are called ``open.''
(One could construct closed universes with flat or hyperboloid geometries by making
topological identifications, but it does not appear to be natural to do so.) Thus, an
intriguing question raised by general relativity is whether our universe is closed or
open. We shall discuss the evidence on this issue at the end of Section~5.4.

Since the isotropic observers are orthogonal to the homogeneous surfaces, we may
express the four-dimensional spacetime metric \(g_{ab}\) as
\begin{equation}
  g_{ab}=-u_a u_b+h_{ab}(t),
  \tag{5.1.10}\label{eq:5.1.10}
\end{equation}
where for each \(t\), \(h_{ab}(t)\) is the metric of either (a) a sphere, (b) flat
Euclidean space, or (c) a hyperboloid, on \(\Sigma_t\). We can choose convenient
coordinates on the four-dimensional spacetime as follows. We choose, respectively,
either (a) spherical coordinates, (b) Cartesian coordinates, or (c) hyperbolic
coordinates on one of the homogeneous hypersurfaces. We then ``carry'' these
coordinates to each of the other homogeneous hypersurfaces by means of our
isotropic observers; i.e., we assign a fixed spatial coordinate label to each observer.
Finally, we label each hypersurface by the proper time, \(\tau\), of a clock carried by
any of the isotropic observers. (By homogeneity, all the isotropic observers must
agree on the time difference between any two hypersurfaces.) Thus, \(\tau\) and our
spatial coordinates label each event in the universe.

Expressed in these coordinates, the spacetime metric takes the form
\begin{equation}
  ds^2=-d\tau^2+a^2(\tau)
  \begin{cases}
    d\psi^2+\sin^2\psi\bigl(d\theta^2+\sin^2\theta\,d\phi^2\bigr),\\
    dx^2+dy^2+dz^2,\\
    d\psi^2+\sinh^2\psi\bigl(d\theta^2+\sin^2\theta\,d\phi^2\bigr),
  \end{cases}
  \tag{5.1.11}\label{eq:5.1.11}
\end{equation}
where the three possibilities in the bracket correspond to the three possible spatial
geometries. [The metric for the spatially flat case could be made to look more
similar to the other cases by writing it in spherical coordinates as
\(d\psi^2+\psi^2(d\theta^2+\sin^2\theta\,d\phi^2)\).] The general form of the
metric, Eq.~\eqref{eq:5.1.11}, is called a Robertson--Walker cosmological model.
Thus, our assumptions of homogeneity and isotropy alone have determined the
spacetime metric up to the three discrete possibilities of spatial geometry and the
arbitrary positive function \(a(\tau)\). To determine the spatial geometry and
\(a(\tau)\), we turn to Einstein's equation.
